\documentclass[14pt,a4paper,oneside]{extbook}
\usepackage[utf8]{vietnam}
\usepackage{amsmath,amssymb,mathrsfs,fancyhdr,enumerate,multirow,makecell}
\usepackage{tkz-euclide,tikz-3dplot,tikz,tkz-tab}
\usetikzlibrary{arrows.meta,calc,intersections,angles,quotes,backgrounds,shapes.geometric,patterns,through,shadows,positioning}
\usepackage{yhmath}
\usepackage[top=0.7cm, bottom=0.7cm, left=1.0cm, right=1.0cm, headheight=16pt, headsep=3mm, footskip=5mm]{geometry}
\usepackage[hidelinks,unicode]{hyperref}
\usepackage{qrcode}
\usepackage[many]{tcolorbox}
\usepackage{fontawesome}

\pagestyle{fancy}
\fancyhf{}
\fancyhead[L]{\footnotesize\bfseries\color{teal!85!black}\faBook~SỔ TAY KIẾN THỨC TOÁN 4}
\fancyhead[R]{\footnotesize\bfseries\color{blue!85!black}Lớp Toán Thầy Lĩnh -- loptoanthaylinh.vn}
\fancyfoot[C]{\footnotesize\bfseries\color{gray!80!black}-- Sổ tay kiến thức: Bài 52. Luyện tập chung thống kê, xác suất --}
\renewcommand{\headrulewidth}{0.4pt}

\begin{document}

% --- BANNER TIÊU ĐỀ BÀI HỌC ---
\noindent
\begin{tcolorbox}[
  enhanced,
  colback=teal!8!white,
  colframe=teal!80!black,
  arc=3mm, boxrule=1pt,
  left=3mm, right=3mm, top=2mm, bottom=2mm,
  halign=center
]
  {\Large\bfseries\color{teal!90!black}\faBookmark\ BÀI 52: LUYỆN TẬP CHUNG THỐNG KÊ, XÁC SUẤT}\\[1.5pt]
  {\footnotesize\itshape\color{black!75} Tổng hợp kĩ năng đọc bảng số liệu, biểu đồ cột và kiểm đếm sự kiện trên 1 trang A4}
\end{tcolorbox}

\vspace{-2.5mm}

\noindent
\begin{minipage}[t]{0.49\textwidth}\vspace{0pt}
\begin{tcolorbox}[
  enhanced,
  colback=blue!4!white,
  colframe=blue!75!black,
  arc=2mm, boxrule=0.8pt,
  left=2mm, right=2mm, top=1.8mm, bottom=1.8mm,
  title={\footnotesize\bfseries\color{white}\faExchange\ 1. Chuyển đổi dữ liệu đa chiều}
]
{\footnotesize
Dữ liệu thống kê có thể chuyển hóa linh hoạt:
\begin{itemize}\setlength{\itemsep}{1pt}\setlength{\leftskip}{-2.5mm}
  \item \textbf{Dãy số liệu} $\to$ Liệt kê các giá trị theo thứ tự thời gian hoặc khảo sát.
  \item \textbf{Bảng số liệu} $\to$ Phân loại thành các hàng và cột khoa học.
  \item \textbf{Biểu đồ cột} $\to$ Trực quan hóa bằng độ cao các cột, giúp so sánh ngay lập tức.
\end{itemize}
Từ bảng số liệu ta có thể dựng biểu đồ cột tương ứng và ngược lại.
}
\end{tcolorbox}
\end{minipage}
\hfill
\begin{minipage}[t]{0.49\textwidth}\vspace{0pt}
\begin{tcolorbox}[
  enhanced,
  colback=teal!4!white,
  colframe=teal!75!black,
  arc=2mm, boxrule=0.8pt,
  left=2mm, right=2mm, top=1.8mm, bottom=1.8mm,
  title={\footnotesize\bfseries\color{white}\faCalculator\ 2. Phương pháp phân tích câu hỏi}
]
{\footnotesize
Bộ 4 câu hỏi phân tích thống kê cốt lõi:
\begin{itemize}\setlength{\itemsep}{1pt}\setlength{\leftskip}{-2.5mm}
  \item \textbf{1. Tìm Max / Min:} Cột cao nhất là Max, cột thấp nhất là Min.
  \item \textbf{2. So sánh hơn kém:} Lấy $\text{Giá trị } A - \text{Giá trị } B$.
  \item \textbf{3. Tính Tổng toàn bộ:} Cộng tất cả số liệu của các cột lại với nhau.
  \item \textbf{4. Tính Trung bình cộng:}
  $$\text{TBC} = \dfrac{\text{Tổng số liệu}}{\text{Số lượng đối tượng}}$$
\end{itemize}
}
\end{tcolorbox}
\end{minipage}

\vspace{-2.5mm}
% --- KHỐI 3: SƠ ĐỒ TƯ DUY TIKZ ---
\noindent
\begin{tcolorbox}[
  enhanced,
  colback=orange!4!white,
  colframe=orange!80!black,
  arc=2mm, boxrule=0.8pt,
  left=2mm, right=2mm, top=1.5mm, bottom=1.5mm,
  title={\footnotesize\bfseries\color{white}\faSitemap\ SƠ ĐỒ TƯ DUY: TAM GIÁC XỬ LÍ DỮ LIỆU THỐNG KÊ \& XÁC SUẤT}
]
\centering
\begin{tikzpicture}[
  node distance=2.2cm,
  block/.style={draw, rounded corners=2pt, align=center, font=\scriptsize\bfseries, inner sep=3pt, minimum width=3.8cm, minimum height=0.9cm},
  arr/.style={<->, {Stealth[length=2mm]}-{Stealth[length=2mm]}, thick}
]
  \node[block, fill=blue!10, draw=blue!80!black] (day) at (0, 1.6) {DÃY SỐ LIỆU\\(Thu thập ban đầu)};
  \node[block, fill=teal!10, draw=teal!80!black] (bang) at (8, 1.6) {BẢNG THỐNG KÊ\\(Phân loại Hàng \& Cột)};
  \node[block, fill=purple!10, draw=purple!80!black] (bieudo) at (4, -0.4) {BIỂU ĐỒ CỘT\\(Trực quan hóa độ cao)};

  \draw[arr, blue!80!black] (day) -- (bang) node[midway, above, font=\tiny\bfseries] {Sắp xếp};
  \draw[arr, teal!80!black] (bang) -- (bieudo) node[midway, right, font=\tiny\bfseries] {Vẽ cột};
  \draw[arr, purple!80!black] (bieudo) -- (day) node[midway, left, font=\tiny\bfseries] {Trích xuất};

  \node[draw=orange!80!black, fill=orange!15, circle, font=\tiny\bfseries, align=center, inner sep=2pt] at (4, 0.9) {KẾT LUẬN\\Min/Max/TBC};
\end{tikzpicture}
\end{tcolorbox}

\vspace{-2.5mm}
% --- KHỐI 4: VÍ DỤ & MẸO TRÁNH BẪY ---
\noindent
\begin{minipage}[t]{0.49\textwidth}\vspace{0pt}
\begin{tcolorbox}[
  enhanced,
  colback=green!4!white,
  colframe=green!60!black,
  arc=2mm, boxrule=0.8pt,
  left=2mm, right=2mm, top=1.8mm, bottom=1.8mm,
  title={\footnotesize\bfseries\color{white}\faCheckCircle\ Ví dụ thực hành chuẩn mực}
]
{\footnotesize
\textbf{Ví dụ:} Số điểm 10 của 4 tổ: Tổ 1: 12 điểm; Tổ 2: 15 điểm; Tổ 3: 18 điểm; Tổ 4: 15 điểm.
\begin{itemize}\setlength{\itemsep}{1pt}\setlength{\leftskip}{-2.5mm}
  \item Tổ đạt nhiều điểm 10 nhất: \textbf{Tổ 3} (18 điểm).
  \item Hai tổ có số điểm bằng nhau: \textbf{Tổ 2 và Tổ 4}.
  \item Cả lớp có tất cả: $12 + 15 + 18 + 15 = \mathbf{60}$ điểm 10.
  \item Trung bình mỗi tổ đạt: $60 : 4 = \mathbf{15}$ điểm 10.
\end{itemize}
}
\end{tcolorbox}
\end{minipage}
\hfill
\begin{minipage}[t]{0.49\textwidth}\vspace{0pt}
\begin{tcolorbox}[
  enhanced,
  colback=red!4!white,
  colframe=red!70!black,
  arc=2mm, boxrule=0.8pt,
  left=2mm, right=2mm, top=1.8mm, bottom=1.8mm,
  title={\footnotesize\bfseries\color{white}\faExclamationTriangle\ Mẹo tránh bẫy trắc nghiệm}
]
{\footnotesize
\begin{itemize}\setlength{\itemsep}{1.5pt}\setlength{\leftskip}{-2.5mm}
  \item \textbf{Bẫy nhầm chia số phần khi tính TBC:}\\
  Khi tính TBC của biểu đồ, phải chia cho đúng \textbf{số cột} (số đối tượng), không được chia nhầm cho tổng số liệu.
  \item \textbf{Bẫy câu hỏi "nhiều hơn/ít hơn":}\\
  Đọc kĩ đề hỏi "nhiều hơn bao nhiêu" (phép trừ) hay hỏi "gấp mấy lần" (phép chia).
\end{itemize}
}
\end{tcolorbox}
\end{minipage}

\end{document}
